% =====================================================================
% 製本用 修士学位論文 表紙（相関基礎科学系）
%
% 出典: 「相関基礎科学系製本用修士学位論文作成要領」
%       ~/Desktop/thesis/master thesis/製本用修士論文作成要領.pdf
%
% 要領の指定:
%   (1) 原則として A4 版。日本語・欧文いずれもパソコンを使用
%   (2) 表紙には次を記載する
%       - 論文題目（題目が外国語の場合は ( ) 書きで日本語訳を必ずつける）
%       - 専攻・系名
%       - 学生証番号
%       - 氏名
%
% 使い方: thesis.tex の \begin{document} 直後で
%     % =====================================================================
% 製本用 修士学位論文 表紙（相関基礎科学系）
%
% 出典: 「相関基礎科学系製本用修士学位論文作成要領」
%       ~/Desktop/thesis/master thesis/製本用修士論文作成要領.pdf
%
% 要領の指定:
%   (1) 原則として A4 版。日本語・欧文いずれもパソコンを使用
%   (2) 表紙には次を記載する
%       - 論文題目（題目が外国語の場合は ( ) 書きで日本語訳を必ずつける）
%       - 専攻・系名
%       - 学生証番号
%       - 氏名
%
% 使い方: thesis.tex の \begin{document} 直後で
%     % =====================================================================
% 製本用 修士学位論文 表紙（相関基礎科学系）
%
% 出典: 「相関基礎科学系製本用修士学位論文作成要領」
%       ~/Desktop/thesis/master thesis/製本用修士論文作成要領.pdf
%
% 要領の指定:
%   (1) 原則として A4 版。日本語・欧文いずれもパソコンを使用
%   (2) 表紙には次を記載する
%       - 論文題目（題目が外国語の場合は ( ) 書きで日本語訳を必ずつける）
%       - 専攻・系名
%       - 学生証番号
%       - 氏名
%
% 使い方: thesis.tex の \begin{document} 直後で
%     % =====================================================================
% 製本用 修士学位論文 表紙（相関基礎科学系）
%
% 出典: 「相関基礎科学系製本用修士学位論文作成要領」
%       ~/Desktop/thesis/master thesis/製本用修士論文作成要領.pdf
%
% 要領の指定:
%   (1) 原則として A4 版。日本語・欧文いずれもパソコンを使用
%   (2) 表紙には次を記載する
%       - 論文題目（題目が外国語の場合は ( ) 書きで日本語訳を必ずつける）
%       - 専攻・系名
%       - 学生証番号
%       - 氏名
%
% 使い方: thesis.tex の \begin{document} 直後で
%     \input{cover_bound}
%     \clearpage
% とする（既存の \maketitle 版の表紙とは別物。審査用と製本用で使い分ける）
%
% ★ 提出前に確認すること
%   - \thesisTitleJa に実際の論文題目を入れる（現在プレースホルダ）
%   - 題目を英語にする場合は \thesisTitleEn を使い、日本語訳を ( ) で必ず併記
%   - \studentID が正しいか学生証で確認する（下の値は過去のファイル名からの推定）
%   - \academicYear は提出年度。2027年3月修了なら「2026年度」
% =====================================================================

\newcommand{\thesisTitleJa}{ここに論文題目を日本語で書く}
% 題目を外国語にする場合はこちらを使い、日本語訳を必ず () で添える:
% \newcommand{\thesisTitleEn}{English Thesis Title}

\newcommand{\academicYear}{2026 年度}
\newcommand{\departmentLine}{広域科学専攻\quad 相関基礎科学系}
\newcommand{\studentID}{31256921}
\newcommand{\authorNameJa}{北岸\quad 健吾}

\begin{titlepage}
\thispagestyle{empty}
\centering
\vspace*{4cm}

{\Large 修士学位論文}

\vspace{3cm}

% --- 論文題目 ---
{\LARGE \thesisTitleJa}

% 題目が外国語のときは上をコメントアウトして次の2行を使う:
% {\LARGE \thesisTitleEn}\\[0.6cm]
% {\large （\thesisTitleJa）}

\vfill

% --- 提出年度・専攻/系・学生証番号・氏名 ---
{\large \academicYear}

\vspace{0.6cm}
{\large \departmentLine}

\vspace{0.6cm}
{\large \studentID}

\vspace{0.6cm}
{\large \authorNameJa}

\vspace{3cm}
\end{titlepage}

%     \clearpage
% とする（既存の \maketitle 版の表紙とは別物。審査用と製本用で使い分ける）
%
% ★ 提出前に確認すること
%   - \thesisTitleJa に実際の論文題目を入れる（現在プレースホルダ）
%   - 題目を英語にする場合は \thesisTitleEn を使い、日本語訳を ( ) で必ず併記
%   - \studentID が正しいか学生証で確認する（下の値は過去のファイル名からの推定）
%   - \academicYear は提出年度。2027年3月修了なら「2026年度」
% =====================================================================

\newcommand{\thesisTitleJa}{ここに論文題目を日本語で書く}
% 題目を外国語にする場合はこちらを使い、日本語訳を必ず () で添える:
% \newcommand{\thesisTitleEn}{English Thesis Title}

\newcommand{\academicYear}{2026 年度}
\newcommand{\departmentLine}{広域科学専攻\quad 相関基礎科学系}
\newcommand{\studentID}{31256921}
\newcommand{\authorNameJa}{北岸\quad 健吾}

\begin{titlepage}
\thispagestyle{empty}
\centering
\vspace*{4cm}

{\Large 修士学位論文}

\vspace{3cm}

% --- 論文題目 ---
{\LARGE \thesisTitleJa}

% 題目が外国語のときは上をコメントアウトして次の2行を使う:
% {\LARGE \thesisTitleEn}\\[0.6cm]
% {\large （\thesisTitleJa）}

\vfill

% --- 提出年度・専攻/系・学生証番号・氏名 ---
{\large \academicYear}

\vspace{0.6cm}
{\large \departmentLine}

\vspace{0.6cm}
{\large \studentID}

\vspace{0.6cm}
{\large \authorNameJa}

\vspace{3cm}
\end{titlepage}

%     \clearpage
% とする（既存の \maketitle 版の表紙とは別物。審査用と製本用で使い分ける）
%
% ★ 提出前に確認すること
%   - \thesisTitleJa に実際の論文題目を入れる（現在プレースホルダ）
%   - 題目を英語にする場合は \thesisTitleEn を使い、日本語訳を ( ) で必ず併記
%   - \studentID が正しいか学生証で確認する（下の値は過去のファイル名からの推定）
%   - \academicYear は提出年度。2027年3月修了なら「2026年度」
% =====================================================================

\newcommand{\thesisTitleJa}{ここに論文題目を日本語で書く}
% 題目を外国語にする場合はこちらを使い、日本語訳を必ず () で添える:
% \newcommand{\thesisTitleEn}{English Thesis Title}

\newcommand{\academicYear}{2026 年度}
\newcommand{\departmentLine}{広域科学専攻\quad 相関基礎科学系}
\newcommand{\studentID}{31256921}
\newcommand{\authorNameJa}{北岸\quad 健吾}

\begin{titlepage}
\thispagestyle{empty}
\centering
\vspace*{4cm}

{\Large 修士学位論文}

\vspace{3cm}

% --- 論文題目 ---
{\LARGE \thesisTitleJa}

% 題目が外国語のときは上をコメントアウトして次の2行を使う:
% {\LARGE \thesisTitleEn}\\[0.6cm]
% {\large （\thesisTitleJa）}

\vfill

% --- 提出年度・専攻/系・学生証番号・氏名 ---
{\large \academicYear}

\vspace{0.6cm}
{\large \departmentLine}

\vspace{0.6cm}
{\large \studentID}

\vspace{0.6cm}
{\large \authorNameJa}

\vspace{3cm}
\end{titlepage}

%     \clearpage
% とする（既存の \maketitle 版の表紙とは別物。審査用と製本用で使い分ける）
%
% ★ 提出前に確認すること
%   - \thesisTitleJa に実際の論文題目を入れる（現在プレースホルダ）
%   - 題目を英語にする場合は \thesisTitleEn を使い、日本語訳を ( ) で必ず併記
%   - \studentID が正しいか学生証で確認する（下の値は過去のファイル名からの推定）
%   - \academicYear は提出年度。2027年3月修了なら「2026年度」
% =====================================================================

\newcommand{\thesisTitleJa}{ここに論文題目を日本語で書く}
% 題目を外国語にする場合はこちらを使い、日本語訳を必ず () で添える:
% \newcommand{\thesisTitleEn}{English Thesis Title}

\newcommand{\academicYear}{2026 年度}
\newcommand{\departmentLine}{広域科学専攻\quad 相関基礎科学系}
\newcommand{\studentID}{31256921}
\newcommand{\authorNameJa}{北岸\quad 健吾}

\begin{titlepage}
\thispagestyle{empty}
\centering
\vspace*{4cm}

{\Large 修士学位論文}

\vspace{3cm}

% --- 論文題目 ---
{\LARGE \thesisTitleJa}

% 題目が外国語のときは上をコメントアウトして次の2行を使う:
% {\LARGE \thesisTitleEn}\\[0.6cm]
% {\large （\thesisTitleJa）}

\vfill

% --- 提出年度・専攻/系・学生証番号・氏名 ---
{\large \academicYear}

\vspace{0.6cm}
{\large \departmentLine}

\vspace{0.6cm}
{\large \studentID}

\vspace{0.6cm}
{\large \authorNameJa}

\vspace{3cm}
\end{titlepage}
